\documentclass[11pt]{article}
\input{packages.tex}
\usepackage{ulem}
\DeclareMathOperator{\sign}{sign}
\DeclareMathOperator{\Imm}{Im}
\DeclareMathOperator{\DIF}{DIF}
\DeclareMathOperator{\rg}{rg}
\DeclareMathOperator{\dist}{dist}

% Настройка колонтитулов
\pagestyle{fancy}
\fancyhf{} % Очистить все поля
\fancyhead[L]{Что ты делаешь в моем болоте? }
\fancyhead[R]{\href{ https://t.me/mallina_olga }{Автор}}
\fancyfoot[C]{\thepage}               % Справа внизу

\setlength{\headheight}{13.6pt}

\hypersetup{
    colorlinks=true,
    linkcolor=blue,
    citecolor=blue,
    filecolor=blue,
    urlcolor=blue,
}

\begin{document}

\begin{titlepage}
    \centering
    \vspace*{\fill}  % заполняет пространство сверху

    {\Huge Лекция 11. Сопряжённые преобразования. Самосопряжённые преобразования. Ортогональные преобразования.
    \par}

    \vspace*{\fill}  % заполняет пространство снизу
\end{titlepage}

\setcounter{page}{1}

\section{Продолжение сопряжённых преобразований из предыдущей серии}
В прошлой серии мы получили следующую формулу:
\[ 1) A^* = \Gamma^{-1} A^T \Gamma \]
Важным частным случаем у нас будет ОНБ:
\[ A^* = A^T \]
\[ 2) \Rightarrow (A^*)^* = A \]
При этом понятно, что  это равенство достаточно проверить в ОНБ, так как это свойство не зависит от базиса.
\[ 3) (AB)^* = B^* A^* \]
\[ 4) det A = det A^T, E=E^T \Rightarrow \text{ у } \varphi \text{ и у }\varphi^* \text{ характеристические многочлены совпадают} \]

\begin{theorem}
    \[ Im\,A = (Ker\,A^*) ^\perp \]
\end{theorem}

\begin{proof}
    Возьмём $\forall x\in E, \forall y\in Ker\,A^*$ и запишем определение сопряжённого преобразования:
    \[ (A(x),y) = (x, A^*(y)) = (x,0) = 0 \]
    Это означает, что
    \[ A(x) \perp \forall y\in Ker\,A^* \Rightarrow Im\, A \subset (Ker\,A^*)^\perp \]
    Пусть $rg A = r$. Найдём $dim\, Ker\, A^*$. Необходимо решить систему:
    \[ A^T y =0  \]
    \[ dim \, Ker\, A^* = n -r \]
    \[ \Rightarrow dim\, (Ker\, A^*)^\perp = n - (n-r) = r \]   
    \[ \Rightarrow dim\, Im\, A = r \]
    Таким образом, мы доказали, что $ Im\, A \subset (Ker\,A^*)^\perp$, но при этом у них равны размерности, а значит, эти подпространства совпадают.
\end{proof}

\section{Самосопряжённые преобразования}
 
\begin{definition}
    Преобразование называется самосопряжённым, если 
    \[ A = A^* \]
    или 
    \[ \forall x, y\in E \ (A(x),y) = (x,A(y)) \]
\end{definition}

\begin{theorem}
    (Критерий самосопряжённости)$A$ --- самосопряжённое $\Leftrightarrow$ его матрица в ОНБ симметрична, то есть $A = A^T$. 
\end{theorem}

\begin{remark}
    Из этой теоремы следует, что для каждого линейного преобразования найдётся ровно одно сопряжённое, те из них, которые равны, будут самосопряжёнными, а те, которые не будут равны --- не будут самосопряжёнными.
\end{remark}

\begin{theorem}
    Все корни характеристического многочлена самосопряжённого преобразования вещественные. 
\end{theorem}

\begin{proof}
    От противного: пусть не все корни вещественные. Тогда $ \exists $ пара комплексных сопряжённых корней. 
    Этой паре соответствует двумерное, инвариантное и не содержащее собственных векторов подпространство $E_1$.
    Рассмотрим ограничение самосопряжённого преобразования на $E_1$. 
    Ограничение --- тоже самосопряжённое. Введём в этом двумерном пространстве ОНБ 
    \[ A' = 
    \begin{pmatrix}
        \alpha \quad \beta \\
        \beta \quad \gamma
    \end{pmatrix} \]
    Тогда характеристический многочлен будет выглядеть так:
    \[ p(\lambda) = (\alpha - \lambda)(\gamma-\lambda) - \beta^2 \]
    \[ p(\alpha) \le 0 \text{, а ветви направлены вверх} \] 
    $\Rightarrow$ есть вещественные корни $\Rightarrow$ есть собственный вектор в этом подпространстве $\Rightarrow$ противоречие.
\end{proof}

\begin{remark}
    Другая формулировка теоремы выше фактически такая: Если $A$ --- симметрична, то уравнение $det(A-\lambda E)$ имеет только вещественные $\lambda$ --- решения.
\end{remark}

\begin{theorem}
    У самосопряжённого преобразования собственные подпространства попарно ортогональны 
    ($\Leftrightarrow$ собственные векторы, отвечающие различным собственным значениям попарно ортогональны )
\end{theorem}

\begin{proof}
    Возьмём два вектора : вектор $h_1$ соответствует собственному значению $\lambda_1$,\\
    $h_2\to \lambda_2, \lambda_1\neq \lambda_2$
    По определению: 
    \[ Ah_1 = \lambda_1 h_1 \]
    \[ Ah_2 = \lambda_2 h_2 \]
    Рассмотрим следующее скалярное произведение: 
    \[ (Ah_1,h_2) = (h_1,Ah_2) \]   
    Но с другой стороны:
    \[ (Ah_1,h_2) = (\lambda_1h_1, h_2) \] 
    \[ (h_1,Ah_2) = (h_1, \lambda_2h_2) \]
    Вычтем из одного другое и получим:
    \[ (\lambda_1 - \lambda_2) (h_1,h_2) = 0 \]
    Но так как $\lambda_1\neq \lambda_2 \Rightarrow (\lambda_1 - \lambda_2) \neq 0$, а значит
    \[ (h_1,h_2) = 0 \Rightarrow h_1\perp h_2 \]
\end{proof}

\begin{theorem}
    Пусть $E_1$ --- инвариантное подпространство самосопряжённого преобразования $A$. Тогда
    $E_1^\perp$ --- тоже инвариантное подпространство $A$. 
\end{theorem}

\begin{proof}
    $E_1$ --- инвариантно 
    \[ \Rightarrow \forall x\in E_1 \ A(x)\in E_1 \]
    Если $\forall y\in E_1^\perp$, то рассмотрим
    \[ 0 = (A(x),y) = (x, A(y)) \Rightarrow A(y)\in E_1^\perp \]   
    $\Rightarrow E_1^\perp$ --- инвариантно.
\end{proof}


\begin{theorem}
    Линейное преобразование $A$ --- самосопряжено $\Leftrightarrow$ в $E \ \exists$ ОНБ из собственных векторов $A$.  
\end{theorem}

\begin{proof}
    $\Leftarrow$:  Если $\exists$ ОНБ из собственных векторов $\to$ матрица преобразования диагональная 
    \[ \begin{pmatrix}
        \lambda_1 &        &        &   \\
        & \ddots &        &   \\
        &        & \ddots &   \\
        &        &        & \lambda_n
    \end{pmatrix} \]
    В ОНБ диагональная --- симметричная $\Rightarrow$ преобразование самосопряжённое \\
    $\Rightarrow$: необходимо доказать, что $E$ = сумма собственных подпространств (как только докажем, в каждом возьмём ОНБ). От противного: пусть $\Sigma$ собственных подпространств = $L \neq E$.
    Тогда $L$ --- инвариантно, а значит, $L^\perp$ --- тоже инвариантно. Рассмотрим ограничение преобразования на $L^\perp$.
    То есть, мы имеем самосопряжённое преобразование на $L^\perp$. Пусть $A'$ --- самосопряжённое преобразование на $L^\perp$. Характеристический многочлен $A'$ имеет только вещественные корни
    $\Rightarrow$ есть собственный вектор, но они все лежат в $L$ --- противоречие. 
\end{proof}

\begin{remark}
    (Эквивалентная формулировка предыдущей теоремы) $\forall$ симметричной матрицы $A \ \exists $ ортогональная матрица $S: S^{-1}AS$ --- диагональная. 
\end{remark}

\section{Ортогональные преобразования}

\begin{definition}
    Преобразование $\varphi$ --- ортогонально, если
    \[ \forall x,y\in E \ (\varphi(x),\varphi(y)) = (x,y) \]
    Свойства этого преобразования: 
    \[ \text{если } x=y \ (\varphi(x),\varphi(x)) = (x,x) \Rightarrow \text{биекция} \]
\end{definition}



\begin{theorem}
    \[ (A(x),A(y)) = (x,y) = (x,A^*Ay) \]
    $\Rightarrow$ преобразование ортогонально $\Leftrightarrow$ $A^*A =E \Leftrightarrow$ в ОНБ $A^* = A^T, A^TA=E$.\\
    То есть преобразование является ортогональным $\Leftrightarrow$ в ОНБ его матрица ортогональна.
\end{theorem}

\begin{theorem}
    $\forall e,f$ ОНБ  $\exists! $ ортогональное преобразование $A$: 
    \[ A(e_i) = f_i \quad \forall i= 1,...,n \]
\end{theorem}

\begin{proof}
    Необходимо посмотреть в ОНБ (e) на следующую матрицу:
    \[ A= \begin{pmatrix}
        &\\
        f_1, ..., f_n\\
        &
    \end{pmatrix} \]
\end{proof}

\begin{theorem}
    Если $A$ --- ортогонально, $\lambda\in \mathbb{R}$ --- собственное значение, тогда $\lambda^2 =1 $.
\end{theorem}

\begin{proof}
    Пусть $h$ --- собственный вектор, отвечающий $\lambda$. Тогда 
    \[ (A(h), Ah) = (\lambda h, \lambda h) = \lambda^2 (h,h)    \]
    С другой стороны
    \[ (A(h), Ah) = (h,h) \]
    \[ \Rightarrow (\lambda^2 -1 ) (h,h) =0 \]
    Но $(h,h) \neq 0 \Rightarrow (\lambda^2 -1 ) =0$
\end{proof}

\begin{theorem}
    Пусть $E_1$ --- инвариантно относительно ортогонального преобразования $A$. Тогда $E_1^\perp$ --- тоже инвариантно относительно $A$. 
\end{theorem}

\begin{proof}
    \[ A(E_1) = E_1 \text{ (из биективности)} \]
    $x\in E_1, y\in E_1^\perp$ 
    \[ 0 = (x,y) = (A(x),A(y)) \Rightarrow A(y)\in A(E_1^\perp) \]
    \[ \Rightarrow dim\, A(E_1^\perp) = dim\, E_1^\perp \Rightarrow A(E_1^\perp) = E_1^\perp \]
\end{proof}

\begin{remark}
    Внимательный слушатель, который был на лекции, знает, что теорему доказать до конца не успели. Автор надеется, что сможет найти в себе силы, упорство, силу духа и что-нибудь еще, чтобы дописать доказательство после следующей лекции (если оно будет не слишком большим, или забьёт на это и понадеется, что другой автор найдёт всё в это в себе)
\end{remark}


\begin{theorem}
    (О полярном разложении) $\forall$ линейное преобразование $A$ евклидова пространства $E$ может быть представлено в виде композиции 
    \[ A = QS \text{ , где } Q \text{--- ортогональное, а } S \text{ --- самосопряжённое с неотрицательными } \lambda \]  
\end{theorem}

\begin{proof}
    Рассмотрим вспомогательное преобразование $\forall A \quad A^*A$, которое является самосопряжённым. 
    У этого преобразования рассмотрим ОНБ $e_1,...,e_n$ из его собственных векторов. Упорядочим эти собственные вектора так, чтобы
    \[ \lambda_1 \ge \lambda_2 \ge ... \ge \lambda_n \]  
    \[ (A(e_i), A(e_j)) = (A^*A(e_i), e_j) = \lambda_i (e_i,e_j) =  
    \begin{cases}
        0, i\neq j\\
        \lambda_i, i =  j
    \end{cases} 
    \]  
    Посмотрим подробнее, что происходит при $i=j$:
    \[ A(e_i), A(e_i) = \lambda_i \ge 0 \Rightarrow |A(e_i)| = \sqrt{\lambda_i} \]
    При этом $\sqrt{\lambda_i}$ называются сингулярными числами.
\end{proof}









\end{document}